% titel: Blatt 7, A4: Druck
% skript: Physik I, Kapitel 8.1.1--8.1.2, Seiten 129--133
% stand: Übungsblatt 7, Aufgabe A4

\begin{werkzeugkasten}[Druck und Tiefe]
In einer ruhenden Flüssigkeit konstanter Dichte \(\rho\) im homogenen
Schwerefeld ist der absolute Druck in der Tiefe \(h\)
\begin{equation}
  p(h)=p_0+\rho gh.
\end{equation}
\(p_0\) ist der Druck an der freien Oberfläche. Druck ist lokal und wirkt
senkrecht auf eine benetzte Fläche, unabhängig von deren Orientierung.
Der Schweredruck \(\rho gh\) ist ein Überdruck gegenüber \(p_0\).

\textit{Skript:} Kap. 8.1.1--8.1.2, S. 129--132, Gl. (8.1) und (8.4).
\end{werkzeugkasten}

\begin{werkzeugkasten}[Von lokalem Druck zur Kraft]
Für eine ebene Fläche mit konstanter Normalenrichtung addieren sich die
senkrechten Beiträge \(\mathrm dF=p\,\mathrm dA\). Ist der Druck überall
gleich, vereinfacht sich das zum Produkt aus Druck und Fläche. Ändert er sich
über die Fläche, muss man die Beiträge integrieren oder den flächengemittelten
Druck verwenden. Druck hat die Einheit \(\si{\pascal}\), Kraft die Einheit
\(\si{\newton}\).

Für die resultierende Belastung eines Wandstücks müssen die Drücke auf beiden
Seiten berücksichtigt werden. Ist der Aussendruck überall gleich \(p_0\),
trägt bei einer offenen Flüssigkeitsoberfläche nur der Schweredruck zur
Druckdifferenz bei. Die reine Kraft der Flüssigkeit auf die Wand umfasst dagegen
auch den Oberflächendruck. Weitere Kräfte wie Halterungen gehören nicht zur
reinen Druckkraft.

\textit{Skript:} Kap. 8.1.1--8.1.2, S. 129--133. Die Integration von
Druckbeiträgen ist eine allgemeine Ergänzung zum dort eingeführten Druckbegriff.
\end{werkzeugkasten}

\begin{achtung}[Darauf achten]
Die Tiefe ist eine vertikale Strecke unter der freien Oberfläche. Für Richtung
und Betrag der Kraft sind die Lage der Fläche und ihre benetzte Seite wichtig;
Fläche und Druck sind nicht dieselbe Grösse. Bei unterschiedlichen Drücken
innen und aussen darf der Luftdruck nicht pauschal weggelassen werden.
\end{achtung}
